\section{Human-Curated SFT Results：Evaluator 会改变结论}

现在回到 human-written Surge 数据实验。课程依次改变 dataset source、prompt/response length 与 dataset size，并同时观察 Open LLM metrics 和 MT-Bench。最重要的结果不是某条柱子最高，而是评价器经常给出不同排序。

\subsection{Open LLM Metrics 与 MT-Bench 的冲突}

上一章已经说明不同 evaluator 回答不同问题，本节用同一批 LLaMA-2-13B checkpoints 检验这种差异是否真的会改变 dataset ranking。读两张柱状图时先固定模型规模与训练阶段，再比较只替换 evaluator 后 apparent winner 是否变化。Open LLM Leaderboard 偏向标准 task metrics，MT-Bench 偏向 judge 对多轮回答的评分；二者权重不同，模型排名自然可能翻转。

\lecturefigure{slide-49.jpg}{不同 SFT datasets 在 Open LLM Leaderboard 上的结果}{官方幻灯片 p.49}

\begin{knowledgebox}{读图：先看每个 benchmark，再看 aggregate}
柱状图由多个 task score 组成，aggregate 高可能由某一强项驱动。先比较 TruthfulQA、reasoning/knowledge 等子项，再检查数据集 task mix；图显示相关性，不单独证明某种数据源导致能力提升。
\end{knowledgebox}

\lecturefigure{slide-50.jpg}{同一批 SFT models 的 MT-Bench scores}{官方幻灯片 p.50}

\begin{warningbox}{Metric reversal 是 measurement signal}
若 Open LLM 与 MT-Bench 选择不同赢家，不应挑一个喜欢的结果。正确做法是定位差异：MT-Bench 是否偏好更长回答？标准 metrics 是否忽略聊天风格？某数据集是否更匹配 judge？冲突本身暴露 evaluation function。
\end{warningbox}

\teachervoice{讲者指出自动 metrics 与 MT-Bench 的结论有时“opposite”。她没有把一个榜单宣布为真值，而是把矛盾带到后面的 length 和 GPT-4 bias 分析。}

\subsection{Response Length 作为 Confounder}

上一小节看到了排名翻转，本节追查一个最容易被忽略的 confounder：不同训练集把模型塑造成了不同回答长度与格式。读长度表时先把它当作 style statistic，而不是质量分，再观察 judge 是否可能奖励这些表面特征。不同 datasets 的平均 completion length 差距很大；若 judge 偏好解释充分、列表化或更长文本，length 就可能同时影响训练 style 与评估分数，形成看似“更 helpful”的假象。

\lecturefigure{slide-51.jpg}{Surge、LIMA 与 OpenAssistant 的平均 response length}{官方幻灯片 p.51}

\begin{knowledgebox}{读图：211、482、722 是 style proxy，不是质量分}
Surge 较短、LIMA 居中、OpenAssistant 较长。数字只描述平均长度；要判断质量，还需控制 task、事实正确率与冗余。把长回答自动当成更全面，会奖励 verbosity 而不是 usefulness。
\end{knowledgebox}

\subsection{Prompt Length 与 Performance}

课程进一步按 average prompt length 观察多个 metrics。长 prompt 可能包含更多上下文，也可能代表更难任务；因此曲线的方向必须和 task composition 一起解释。

\lecturefigure{slide-52.jpg}{Performance 与 average prompt length 的关系}{官方幻灯片 p.52}

\begin{knowledgebox}{读图：多面板曲线先看一致性}
先检查各 benchmark 斜率是否同向，再找 turning point 与离群线。若只有某一 metric 上升，说明 length 可能改善该 metric 所偏好的 behavior，而不是全面能力。曲线也不控制 task difficulty。
\end{knowledgebox}

\lecturefigure{slide-53.jpg}{不同 prompt-length setting 的 MT-Bench scores}{官方幻灯片 p.53}

\begin{warningbox}{LLM judge 对 length 的反应并不稳定}
课堂发现 MT-Bench 与部分自动 metrics 再次不一致，且 GPT-4 未简单地总偏好更长 prompt 对应的模型。Judge bias 是多维的：response verbosity、格式、训练风格和 task 都会共同作用。
\end{warningbox}

\subsection{Dataset Size Ablation 与 Diminishing Returns}

最后一组 human-data 实验从 Surge-Instruct 抽取不同规模子集。Ablation 比跨数据集比较更接近因果，因为 provenance 和 task taxonomy 更一致，但抽样方差、训练 tokens 与 evaluator 仍需考虑。

\lecturefigure{slide-54.jpg}{Surge-Instruct dataset size ablation 的自动指标}{官方幻灯片 p.54}

\begin{knowledgebox}{读图：看边际斜率与置信不确定性}
随着样本增加，多条曲线趋于平台，支持 diminishing returns。图没有显示所有随机种子与置信区间，因此小差值不宜过度解释；更可靠的是“大部分收益早出现、继续加相似数据收益变小”的趋势。
\end{knowledgebox}

\lecturefigure{slide-55.jpg}{相同 dataset-size ablation 的 MT-Bench 结果}{官方幻灯片 p.55}

\begin{knowledgebox}{读图：同一 ablation 也会被 evaluator 重排}
MT-Bench 的变化不必与自动指标同向。若增加数据改善标准 task 却不改善 judge score，可能是 task coverage 与 conversational style 的分离；反过来也可能是 judge style reward。实验结论应写成 metric vector，而不是一个总分。
\end{knowledgebox}

\teachervoice{Rajani 的总结是 TruthfulQA 在自动 metrics 中最能拉开差异，而 MT-Bench 与这些 metrics 相关性有限。她把这个现象视为 benchmark gap，而不是简单宣布 MT-Bench 或传统 benchmark 失败。}

\subsection{本章小结}

Human-curated SFT 的数据源、长度和规模都会影响模型，但 evaluator 会显著改变 apparent winner。高质量实验应同时报告 task metrics、judge scores、response statistics 和 human evidence，并把不一致当作待解释结果。

